\chapter{自我与人格}

\section{自我领域}

自我（Self）维护主体对自身长期身份和特征的认知。自我领域至少包含自我模型、倾向与人格风格、叙事身份三类表示；这些表示属于同一领域，可以共享来源、修订和查询基础。

\begin{description}
\item[自我模型] 表达主体身份、角色、能力、边界和对自身当前状态的长期认识。
\item[倾向与人格风格] 表达相对稳定的偏好、价值倾向、习惯性适应和互动风格。人格风格（Persona）是自我领域中的表示，不建立并列的顶层权威。
\item[叙事身份] 表达主体对自身经历、变化、连续性和长期方向的组织性理解。
\end{description}

\section{认知种子}

认知种子（Cognitive Seed）用于在主体供给时建立初始自我认知和倾向。种子内容必须具有来源和版本，进入自我领域后按照自我修订合同演化。

产品安全规则、操作权限和行为禁令属于宿主或 Heptalogos 的系统宪制，不作为人格或自我认知保存。普通认知变化不能覆盖系统宪制。

本设计不预设独立于普通自我修订机制的“不可演化认知人格核心”。如未来需要引入更强的构成性认知状态，必须通过独立设计决定其身份、修改权限和与系统宪制的边界。

\section{自我与其他领域}

长期稳定的价值倾向、偏好和人格特征由自我领域维护；面向具体未来状态、可被激活或满足的愿望与目标由动机领域维护。具体经历和关于自身的事件仍可以由记忆域保存，并作为自我修订的证据。

社会关系对主体行为的影响由社会认知提供关系上下文，自我领域不复制针对具体对象的关系状态。行为表达由 Heptalogos 根据自我、社会、动机和当前情境共同决定。
